\documentclass[aspectratio=169,11pt]{beamer}
\usepackage[UTF8,fontset=fandol]{ctex}
\usepackage{amsmath,amssymb,mathtools,booktabs}
\usepackage{tikz,pgfplots}
\usetikzlibrary{arrows.meta,calc,positioning}
\pgfplotsset{compat=1.18}
\usepackage{lesson-theme}

\title{拉格朗日中值定理}
\subtitle{定理、证明与应用}
\author{高等数学 · 定理与应用}
\date{}

\begin{document}

% 1
\LessonBackground{assets/cover-bg.png}
\begin{frame}[plain]
  \begin{tikzpicture}[remember picture,overlay]
    \node[anchor=west,text width=.54\paperwidth] at ([xshift=13mm]current page.west) {
      {\usebeamerfont{title}\color{LessonInk}\inserttitle\par}
      \vspace{3mm}
      {\usebeamerfont{subtitle}\color{LessonAccent}\insertsubtitle\par}
      \vspace{9mm}
      {\color{LessonMuted}\insertauthor}
    };
  \end{tikzpicture}
\end{frame}
\LessonClearBackground

% 2
\begin{frame}{课程目标}
  \begin{columns}[T,onlytextwidth]
    \column{.58\textwidth}
      \begin{block}{学完后，你能够}
        \begin{itemize}
          \item 准确复述定理的\alert{条件、结论与区间}
          \item 用辅助函数把问题转化为罗尔定理
          \item 用中值定理证明函数不等式与变化率估计
        \end{itemize}
      \end{block}
    \column{.36\textwidth}
      \textbf{前置知识}
      \vspace{2mm}
      \begin{itemize}
        \item 连续与可导
        \item 导数的几何意义
        \item 罗尔定理
      \end{itemize}
  \end{columns}
  \vfill
  \LessonTakeaway{今天只回答一个核心问题：平均变化率会不会在某一点被瞬时变化率取到？}
\end{frame}

% 3
\begin{frame}{问题引入}
  \begin{columns}[c,onlytextwidth]
    \column{.58\textwidth}
      \centering
      \begin{tikzpicture}
        \begin{axis}[
          width=8.1cm,height=5.15cm,
          xmin=0,xmax=5,ymin=0,ymax=4.5,
          axis lines=middle,ticks=none,
          axis line style={LessonMuted!55},
          clip=false]
          \addplot[LessonAccent,very thick,domain=.35:4.65,samples=80]{.18*x^2+.35};
          \addplot[LessonMuted,thick,domain=.7:4.3]{.9*x-.192};
          \addplot[LessonTeal,very thick,domain=1.45:3.55]{.9*x-.775};
          \addplot[only marks,mark=*,mark size=2.2pt,LessonAccent] coordinates {(.7,.438) (4.3,3.678)};
          \addplot[only marks,mark=*,mark size=2.2pt,LessonTeal] coordinates {(2.5,1.475)};
          \node[below] at (axis cs:.7,.1) {$a$};
          \node[below] at (axis cs:2.5,.1) {$\xi$};
          \node[below] at (axis cs:4.3,.1) {$b$};
          \node[text=LessonMuted] at (axis cs:3.9,3.05) {割线};
          \node[text=LessonTeal] at (axis cs:1.75,1.1) {切线};
        \end{axis}
      \end{tikzpicture}
    \column{.38\textwidth}
      \textbf{区间平均变化率}
      \[
        \frac{f(b)-f(a)}{b-a}
      \]
      \textbf{点处瞬时变化率}
      \[
        f'(\xi)
      \]
      \vspace{2mm}
      \LessonQuestion{在什么条件下，存在 $\xi\in(a,b)$ 使二者相等？}
  \end{columns}
\end{frame}

% 4
\begin{frame}{定理}
  \begin{block}{拉格朗日中值定理}
    设 $a<b$。若函数 $f$ 在闭区间 $[a,b]$ 上\alert{连续}，在开区间 $(a,b)$ 内\alert{可导}，则至少存在一点 $\xi\in(a,b)$，使得
    \[
      \boxed{\;f'(\xi)=\frac{f(b)-f(a)}{b-a}\;}.
    \]
  \end{block}
  \vspace{1mm}
  \begin{center}
    \large 等价地，\qquad
    $\boxed{\;f(b)-f(a)=f'(\xi)(b-a)\;}$
  \end{center}
  \vfill
  \LessonTakeaway{定理保证“至少存在”，但不保证 $\xi$ 唯一，也不保证它是中点。}
\end{frame}

% 5
\begin{frame}{几何意义}
  \begin{columns}[c,onlytextwidth]
    \column{.64\textwidth}
      \centering
      \begin{tikzpicture}
        \begin{axis}[
          width=8.8cm,height=5.2cm,
          xmin=0,xmax=5,ymin=0,ymax=4.35,
          axis lines=none,ticks=none,clip=false]
          \addplot[LessonAccent,very thick,domain=.35:4.65,samples=80]{.18*x^2+.35};
          \addplot[LessonMuted,thick,domain=.7:4.3]{.9*x-.192};
          \addplot[LessonMuted!25,thick,domain=1.0:4.0]{.9*x-.46};
          \addplot[LessonTeal,very thick,domain=1.35:3.65]{.9*x-.775};
          \draw[-{Latex[length=2.2mm]},LessonWarm,thick] (axis cs:4.45,3.55) -- (axis cs:4.45,2.98);
          \draw[-{Latex[length=2.2mm]},LessonWarm,thick] (axis cs:4.45,2.85) -- (axis cs:4.45,2.32);
          \addplot[only marks,mark=*,mark size=2pt,LessonTeal] coordinates {(2.5,1.475)};
          \node[text=LessonTeal,below] at (axis cs:2.5,1.42) {$\xi$};
          \node[text=LessonMuted] at (axis cs:3.7,3.25) {固定斜率};
          \node[text=LessonTeal] at (axis cs:1.8,.72) {相切};
        \end{axis}
      \end{tikzpicture}
    \column{.32\textwidth}
      \begin{enumerate}
        \item 连接端点，锁定割线斜率
        \item 保持方向，平行移动直线
        \item 接触曲线时，切线与割线平行
      \end{enumerate}
      \vspace{4mm}
      \[
        f'(\xi)=\frac{f(b)-f(a)}{b-a}
      \]
  \end{columns}
\end{frame}

% 6
\begin{frame}{辅助函数}
  \begin{columns}[c,onlytextwidth]
    \column{.57\textwidth}
      过 $(a,f(a))$ 与 $(b,f(b))$ 的割线为
      \[
        \ell(x)=f(a)+\frac{f(b)-f(a)}{b-a}(x-a).
      \]
      构造竖直差
      {\small\[
        \boxed{F(x)=f(x)-f(a)-\frac{f(b)-f(a)}{b-a}(x-a)}.
      \]}
      于是 $F(a)=F(b)=0$。
    \column{.39\textwidth}
      \centering
      \begin{tikzpicture}[x=.93cm,y=.80cm]
        \draw[->,LessonMuted!60] (-.2,0) -- (5.15,0) node[right] {$x$};
        \draw[LessonAccent,very thick,domain=.35:4.65,samples=70] plot (\x,{.18*(\x-.7)*(\x-4.3)});
        \draw[LessonTeal,very thick] (1.55,-.583) -- (3.45,-.583);
        \fill[LessonAccent] (.7,0) circle (1.8pt) node[above] {$a$};
        \fill[LessonAccent] (4.3,0) circle (1.8pt) node[above] {$b$};
        \fill[LessonTeal] (2.5,-.583) circle (1.8pt) node[below] {$\xi$};
        \node[text=LessonAccent] at (4.45,-.95) {$F$};
        \node[text=LessonTeal] at (2.5,-.2) {$F'(\xi)=0$};
      \end{tikzpicture}
  \end{columns}
  \vspace{-0.5mm}
  \vfill
  \LessonTakeaway{$F$ 满足罗尔定理的端点条件。}
\end{frame}

% 7
\LessonBackground{assets/middle-bg.png}
\begin{frame}{定理证明}
  \begin{enumerate}
    \item $f$ 与 $\ell$ 在 $[a,b]$ 连续、在 $(a,b)$ 可导，故 $F=f-\ell$ 也满足同样条件。
    \vspace{1.2mm}
    \item 由构造直接得到
      \[
        F(a)=F(b)=0.
      \]
    \item 由罗尔定理，存在 $\xi\in(a,b)$，使
      \[
        F'(\xi)=0.
      \]
    \item 计算导数并整理：
      \[
        0=F'(\xi)=f'(\xi)-\frac{f(b)-f(a)}{b-a}
        \quad\Longrightarrow\quad
        \boxed{f(b)-f(a)=f'(\xi)(b-a)}.
      \]
  \end{enumerate}
\end{frame}
\LessonClearBackground

% 8
\begin{frame}{使用条件}
  \begin{columns}[T,onlytextwidth]
    \column{.48\textwidth}
      \textbf{使用清单}
      \vspace{2mm}
      \begin{itemize}
        \item[$\checkmark$] 先写清 $a<b$
        \item[$\checkmark$] $f$ 在 $[a,b]$ 上连续
        \item[$\checkmark$] $f$ 在 $(a,b)$ 内可导
        \item[$\checkmark$] 结论中的 $\xi$ 位于开区间
      \end{itemize}
      \vspace{2mm}
      \textcolor{LessonTeal}{端点处不要求可导。}
    \column{.48\textwidth}
      \textbf{常见误判}
      \vspace{2mm}
      \begin{itemize}
        \item[\textcolor{LessonBad}{$\times$}] 把 $\xi$ 直接取成 $\frac{a+b}{2}$
        \item[\textcolor{LessonBad}{$\times$}] 只检查可导，不检查闭区间连续
        \item[\textcolor{LessonBad}{$\times$}] 认为条件不满足时结论一定不成立
        \item[\textcolor{LessonBad}{$\times$}] 忘记说明“至少存在一点”
      \end{itemize}
  \end{columns}
  \vfill
  \LessonTakeaway{条件是保证结论成立的充分条件；缺一项时，不能直接套用定理。}
\end{frame}

% 9
\begin{frame}{常用推论}
  \textbf{有限增量公式}\hfill
  若 $x$ 与 $x+\Delta x$ 之间满足定理条件，则存在 $\theta\in(0,1)$：
  \[
    \boxed{f(x+\Delta x)-f(x)=f'(x+\theta\Delta x)\,\Delta x}.
  \]
  \vspace{1mm}
  \begin{columns}[T,onlytextwidth]
    \column{.48\textwidth}
      \textbf{导数有界 $\Rightarrow$ 变化有界}
      \[
        |f'(t)|\le M
        \Longrightarrow
        |f(b)-f(a)|\le M|b-a|.
      \]
    \column{.48\textwidth}
      \textbf{导数符号 $\Rightarrow$ 单调性}
      \[
        f'(t)\ge 0\Rightarrow f\text{ 单调不减},
      \]
      \[
        f'(t)>0\Rightarrow f\text{ 严格递增}.
      \]
  \end{columns}
  \vfill
  \LessonTakeaway{中值定理把局部导数信息转化为整体函数变化。}
\end{frame}

% 10
\begin{frame}{例 1：对数不等式}
  \textbf{证明：}对任意 $x>0$，证明 $\ln x\le x-1$。
  \vspace{0mm}
  \begin{itemize}
    \item 当 $x=1$ 时，等号显然成立。
    \item 当 $x\ne1$ 时，对 $g(t)=\ln t$ 在端点 $1,x$ 之间应用中值定理：
  \end{itemize}
  \[
    \exists\,\xi\in(\min\{1,x\},\max\{1,x\}),\qquad
    \ln x=\frac{x-1}{\xi}.
  \]
  \begin{columns}[T,onlytextwidth]
    \column{.47\textwidth}
      \centering
      $x>1\Rightarrow \xi>1$
      \[
        \frac{x-1}{\xi}<x-1.
      \]
    \column{.47\textwidth}
      \centering
      $0<x<1\Rightarrow \xi<1$，且 $x-1<0$
      \[
        \frac{x-1}{\xi}<x-1.
      \]
  \end{columns}
  \vfill
  \LessonTakeaway{$\ln x<x-1$ 对 $x\ne1$ 成立；等号只在 $x=1$ 取得。}
\end{frame}

% 11
\begin{frame}{例 2：正弦函数估计}
  \textbf{证明：}对任意 $x,y\in\mathbb{R}$，证明
  \[
    |\sin x-\sin y|\le |x-y|.
  \]
  \vspace{0mm}
  若 $x=y$，结论显然。若 $x\ne y$，函数 $f(t)=\sin t$ 在 $x,y$ 之间连续且可导，因此存在
  \[
    \xi\in(\min\{x,y\},\max\{x,y\})
  \]
  使得
  \[
    \sin x-\sin y=\cos\xi\,(x-y).
  \]
  取绝对值：
  \[
    \boxed{|\sin x-\sin y|=|\cos\xi|\,|x-y|\le |x-y|}.
  \]
  \vfill
  \LessonTakeaway{$|f'|\le1$ 意味着输入变化多少，函数值最多变化同样多。}
\end{frame}

% 12
\begin{frame}{课堂练习}
  \begin{enumerate}
    \item $f(x)=|x|$ 在 $[-1,1]$ 上能否应用拉格朗日中值定理？说明理由。
    \vspace{4mm}
    \item 若 $f$ 在 $[-1,2]$ 上连续、在 $(-1,2)$ 内可导，且 $|f'(x)|\le3$，则
      \[
        |f(2)-f(-1)|\le\underline{\hspace{16mm}}.
      \]
    \vspace{3mm}
    \item 判断：定理中的 $\xi$ 总可以取区间中点。\hspace{4mm}
      $\square$ 正确\qquad$\square$ 错误
  \end{enumerate}
  \vfill
  \LessonQuestion{每一题都先说“条件是否满足”，再说“结论如何使用”。}
\end{frame}

% 13
\LessonBackground{assets/closing-bg.png}
\begin{frame}{课程回顾}
  \begin{columns}[c,onlytextwidth]
    \column{.64\textwidth}
      \begin{beamercolorbox}[rounded=true,sep=3.2mm,wd=\linewidth]{block body}
        \textbf{条件}\quad $[a,b]$ 连续，$(a,b)$ 可导
        \[
          \Downarrow\quad \exists\,\xi\in(a,b)\quad\Downarrow
        \]
        \[
          \boxed{f(b)-f(a)=f'(\xi)(b-a)}
        \]
        \vspace{0mm}
        \begin{itemize}
          \item \textbf{理解：}割线斜率 = 某处切线斜率
          \item \textbf{证明：}减去割线，再调用罗尔定理
          \item \textbf{应用：}估计变化、判断单调性、证明不等式
        \end{itemize}
        \vspace{0mm}
        \textcolor{LessonAccent}{\textbf{解题顺序：}定区间 $\rightarrow$ 验条件 $\rightarrow$ 选等式或上界}
      \end{beamercolorbox}
    \column{.32\textwidth}
      \mbox{}
  \end{columns}
\end{frame}
\LessonClearBackground

\end{document}
